\chapter{绪\hspace{6pt}论}

\section{研究工作的背景与意义}

随着互联网与数字化技术的飞速发展，海量关于人物、组织、地点及事件的信息不断积累。然而，这些数据多以非结构化的形式散落于不同来源，难以被机器直接理解与计算。因此，如何将庞杂且孤立的数据转化为可关联、可查询的系统性知识，不仅是突破信息过载瓶颈的关键，更是实现认知智能的必由之路。为应对这一挑战，知识图谱（Knowledge Graph，KG）应运而生。它以实体及其关系为基本单元，通过（头实体，关系，尾实体）三元组的形式精确描述客观事实，从而将碎片化的信息编织成结构化的语义网络。这种表示方式不仅显式刻画了对象的属性及其多维联系，更为机器实现深度推理与大规模数据利用奠定了不可或缺的基础\cite{KGReview_Nickel_2016}。

近年来，学术界与工业界相继构建了多种具有代表性的大规模知识图谱，并已被广泛应用于智能搜索\cite{singhal2012knowledgegraph}、自动问答\cite{berant2013semantic}与个性化推荐\cite{wang2019kgat}等下游任务中。例如，Freebase 通过社区协作，汇集了人物、影视、地理等多个领域的结构化事实\cite{Freebase_Bollacker_2008}。DBpedia 则从 Wikipedia 中抽取结构化信息，并将其与其他开放数据集深度关联，从而有效支持跨数据源的语义查询\cite{auer2007dbpedia}。YAGO 创新性地融合了 Wikipedia 的实体信息与 WordNet 的概念分类体系，建立了实体与抽象语义类别间的精准映射\cite{suchanek2007yago}。此外，Wikidata 采用开放协作模式，维护着一个可供多语言版本及外部应用共享的结构化知识库\cite{vrandecic2014wikidata}。这些知识图谱在数据来源与组织方式上各有侧重，共同推动了全球结构化知识资源的积累与复用。

尽管这些知识图谱已汇聚了海量的结构化事实，但受限于客观世界的开放性与动态演化，现有图谱在知识覆盖上远未达到完备，普遍面临着严重的信息缺失问题。一方面，自动化信息抽取受限于语料的覆盖范围及算法的性能瓶颈，难以将所有客观事实均转化为可靠的三元组记录；另一方面，人工校验与维护不仅面临高昂的成本约束，也难以实时跟进实体属性与关联信息的动态演化。受上述构建瓶颈制约，现有知识库往往只能捕获现实世界的局部切片。以 Freebase 为例，针对具有多个取值的复杂关系，图谱通常仅能收录部分核心对象，难以穷尽完整的事实集合\cite{dong2014knowledgevault}。更值得注意的是，事实缺失在不同知识图谱及不同属性之间呈现出显著的不均衡性，并集中表现为基础属性和关系信息的遗漏。已有统计显示，Freebase 与 DBpedia 中分别约有 71\% 和 66\% 的人物实体缺少出生地信息，DBpedia 中约有 58\% 的科学家缺少研究领域信息\cite{krompass2015dissertation}。此类缺失还不局限于实例层的关系事实，实体类型、关系的定义域与值域等模式层信息同样可能不完整\cite{krompass2015typeconstrained}。多层次的信息缺失及其不均衡分布进一步加剧了知识图谱的结构稀疏性，并限制了依赖关系连接与路径信息的检索、问答和图分析任务\cite{KGReview_Nickel_2016}。

知识图谱推理（Knowledge Graph Reasoning）旨在依据图谱中已有的实体、关系、规则及其他约束信息，推断尚未被显式记录的事实，是缓解知识图谱不完备性的重要技术途径。根据推理机制的不同，现有方法既包括基于本体、公理和逻辑规则的符号推导，也包括借助表示学习、关系路径和神经网络进行的统计推断。从任务层面来看，知识图谱补全与知识图谱问答均需要在已有知识的基础上推断目标事实，是知识图谱推理的两类典型任务\cite{zhang2021neuralreasoning}。

其中，知识图谱补全（Knowledge Graph Completion, KGC）直接面向图谱中的事实缺失问题，通过学习已有事实中的关系规律，并结合可用的辅助信息，对尚未记录的候选事实进行预测与可信度评估\cite{bordes2013translating,KGReview_Nickel_2016}。其核心价值在于恢复实体间的潜在语义联系，拓宽知识覆盖边界，从而为下游任务提供更为坚实的事实依据。具体而言，自动化的预测机制不仅能高效挖掘已有图谱中的潜在关系，大幅降低人工逐项补录与维护的成本；更重要的是，经核验补入的可靠事实能将原本孤立的节点重新连入图网络中，为知识问答、关联检索与个性化推荐等下游应用提供全局连通的结构化支撑。考虑到预测结果中的误差可能沿图结构传播并干扰下游系统，知识图谱补全模型必须在知识拓展的广度与预测的可靠性之间取得严格平衡，而非单纯追求新增三元组的数量规模\cite{dong2014knowledgevault}。

在具体的研究与评测中，知识图谱补全通常划分为三元组分类、关系预测和链接预测三大类任务。三元组分类旨在判断给定三元组事实的真伪；关系预测旨在推断已知实体对之间的缺失关系；而链接预测则是在给定关系及某一端实体的条件下，对缺失的头实体或尾实体进行预测与排序\cite{yao2019kgbert}。相较于前两项任务，链接预测不仅要求模型具备对图谱拓扑结构的深度理解，还直接面向真实知识库构建中最为广泛的“实体补漏”需求，能够提供更全面、细粒度的推理结果。因此，链接预测被广泛视为知识图谱补全领域中最核心、最具代表性的基础任务。

知识图谱嵌入（Knowledge Graph Embedding, KGE）是实现知识图谱补全的经典技术路线。该类方法将实体和关系映射到低维连续向量空间，并利用评分函数刻画三元组的合理性，从而从已有事实中学习潜在的关系模式\cite{bordes2013translating}。KGE 能够有效利用知识图谱中的拓扑结构信息，并对大规模候选实体进行高效的计算与排序。然而，当模型主要以实体和关系标识（ID）作为输入时，其表示质量高度依赖已有图结构的完整度，难以直接利用实体名称、文本描述和关系语义中蕴含的丰富背景知识。预训练语言模型的发展，恰为引入此类文本信息、弥补结构证据的不足提供了全新的技术基础。

近年来，预训练语言模型（Pre-trained Language Model, PLM）在模型架构、训练规模和任务适配方式等方面取得了突破性进展。Transformer 利用自注意力机制深入建模序列内部的依赖关系，为大规模语言表示学习奠定了架构基础\cite{vaswani2017attention}。BERT 通过双向上下文预训练获得了具备强大迁移能力的语义表示，推动了预训练模型在多类语言理解任务中的应用\cite{devlin2019bert}；RoBERTa 则通过扩大训练数据规模并优化预训练策略，进一步提升了模型的文本表示能力\cite{liu2019roberta}。在此基础上，T5 将不同的自然语言处理任务统一为文本到文本（Text-to-Text）的学习范式，增强了同一模型处理理解与生成任务的通用性\cite{raffel2020t5}；GPT-3 则展示了超大规模模型仅需少量上下文示例即可适配多种任务的强大上下文学习能力\cite{brown2020gpt3}。随着参数规模、数据量级和任务泛化能力的持续扩张，预训练语言模型逐步演进为大语言模型（Large Language Model, LLM）。其中，OpenAI 发布的 ChatGPT 以多轮对话形式展现了大语言模型的通用交互与任务处理能力，成为大模型发展的代表性进展\cite{openai2022chatgpt}。大语言模型的作用也由提供通用文本表示，拓展至利用预训练知识、上下文学习和文本生成能力来解决复杂推理任务。

对于知识图谱补全而言，实体名称、描述文本和关系标签能够提供详实的实体类型、属性特征及背景语义，这为结构稀疏或局部关系不足的长尾实体补充了关键的外部证据。基于预训练语言模型的方法既可以对这些文本信息进行深度的语义编码，也可以将缺失事实的预测任务转化为文本判别、语义匹配或条件生成问题。以 KG-BERT 为代表的早期工作验证了这一思路，其通过对三元组文本进行上下文编码来判断候选事实的合理性\cite{yao2019kgbert}。该类研究表明，预训练模型学习到的语言表示能够深度参与结构化事实的预测，这也为后续从语义编码、文本生成和提示引导（Prompting）等角度利用大模型开展 KGC 研究奠定了基础。

然而，语言模型提供的语义信息并不能完全替代知识图谱中的显式结构约束。一方面，实体间描述的文本相似性并不等同于它们在客观世界中具有真实的关系连接，且预训练语料库中蕴含的静态知识未必与目标知识图谱的事实边界及更新状态完全吻合；另一方面，单纯依赖传统图结构又易受事实缺失、长尾分布和实体邻域稀疏等问题的制约。因此，如何在充分发挥预训练语言模型语义表征能力的同时，有效保留并融合知识图谱中严谨的关系模式与拓扑证据，成为提升补全准确性与可靠性的关键挑战。

基于上述背景，本文以知识图谱补全作为知识推理的具体研究切入点，并聚焦于其最核心的链接预测任务。本文将重点研究如何协同建模现有的图拓扑结构（结构证据）与丰富的实体文本信息（文本语义），从而对缺失事实所指向的候选实体进行精准判断与排序。

% TODO（原 1.1.3：本课题的理论价值与实际意义）：待研究内容确定后，
% 在本节末尾补充本课题的理论价值与实际意义，与上述背景自然衔接，不单列三级标题。



\section{时域积分方程方法的国内外研究历史与现状}
时域积分方程方法的研究始于上世纪60 年代，C.L.Bennet 等学者针对导体目
标的瞬态电磁散射问题提出了求解时域积分方程的时间步进（marching-on in-time,
MOT）算法。

\section{本文的主要贡献与创新}
本论文以时域积分方程时间步进算法的数值实现技术、后时稳定性问题以及两层平面波加速算法为重点研究内容，主要创新点与贡献如下：

\section{本论文的结构安排}
本文的章节结构安排如下：